\section*{Capítulo \ref{ch:b2:blood-circulation} --- La sangre y el sistema circulatorio}

\begin{solution}{exo:b2:blood-circulation:1}
\omterm{def:b2:blood-circulation:blood}{Plasma} \qty{55}{\%} (agua, \qty{70}{g/L} de proteínas, sales, nutrientes,
gases, hormonas); células \qty{45}{\%}. \omterm{def:b2:blood-circulation:blood}{Glóbulos rojos}, $5\times
10^{12}$/L, discos de \qty{7}{\micro m}, transporte de oxígeno; glóbulos
blancos, $7\times 10^{9}$/L, \qtyrange{10}{15}{\micro m}, inmunidad;
\omterm{def:b2:blood-circulation:blood}{plaquetas}, $3\times 10^{11}$/L, fragmentos de \qty{2}{\micro m},
hemostasia.
\end{solution}

\begin{solution}{exo:b2:blood-circulation:2}
Pez: corazón $\to$ branquias $\to$ cuerpo $\to$ corazón, un solo circuito,
en el que el cuerpo recibe la \omterm{def:b2:blood-circulation:blood}{sangre} con la baja presión que queda tras
las branquias. Mamífero: corazón derecho $\to$ pulmones (\qty{25}{mmHg})
$\to$ corazón izquierdo $\to$ cuerpo
(\qty{120}{mmHg}) $\to$ corazón derecho. El circuito doble da al cuerpo una
presión alta sin exponer a ella los pulmones, y permite regular los dos
circuitos por separado.
\end{solution}

\begin{solution}{exo:b2:blood-circulation:3}
\omterm{def:b2:blood-circulation:vessels}{Arteria}: pared gruesa con capas elásticas y músculo --- conduce la \omterm{def:b2:blood-circulation:blood}{sangre}
a alta presión y suaviza el pulso. \omterm{def:b2:blood-circulation:vessels}{Arteriola}: pared sobre todo de músculo
liso --- fija la resistencia y distribuye el flujo. \omterm{def:b2:blood-circulation:vessels}{Capilar}: una sola
célula endotelial de grosor --- intercambio. \omterm{def:b2:blood-circulation:vessels}{Vena}: pared delgada y
distensible, con válvulas --- devuelve la \omterm{def:b2:blood-circulation:blood}{sangre} a baja presión y almacena
la mayor parte de ella.
\end{solution}

\begin{solution}{exo:b2:blood-circulation:4}
$70\times 72 = \qty{5}{L/min}$, \qty{7200}{L} al día, frente a \qty{5}{L}
de \omterm{def:b2:blood-circulation:blood}{sangre}: el volumen pasa por el corazón \num{1440} veces al día. Ningún
órgano podría fabricar ni consumir siete toneladas de \omterm{def:b2:blood-circulation:blood}{sangre} al día; la
misma \omterm{def:b2:blood-circulation:blood}{sangre} debe volver --- circula.
\end{solution}

\begin{solution}{exo:b2:blood-circulation:5}
$\Delta P = 8\eta LQ/\pi r^{4} = 8\times 3\times 10^{-3}\times
0.4\times 8.3\times 10^{-5}/(\pi\times 2.44\times 10^{-8}) =
\qty{10}{Pa}$, \qty{0.08}{mmHg}: la aorta no cuesta nada; la presión se
gasta en las \omterm{def:b2:blood-circulation:vessels}{arteriolas}.
\end{solution}

\begin{solution}{exo:b2:blood-circulation:6}
$(15/12)^{4} = 2.4$: la resistencia se multiplica por 2,4 y el caudal baja
al \qty{41}{\%}. $(15/20)^{4} = 0.32$: la resistencia baja a un tercio y
el caudal se multiplica por 3,2.
\end{solution}

\begin{solution}{exo:b2:blood-circulation:7}
Aorta: $83/3 = \qty{28}{cm/s}$; \omterm{def:b2:blood-circulation:vessels}{capilares}: $83/3000 =
\qty{0.028}{cm/s}$; tránsito: $0.1/0.028 = \qty{3.6}{s}$. Durante el
ejercicio el gasto se multiplica por cinco y la superficie \omterm{def:b2:blood-circulation:vessels}{capilar} abierta
quizá por tres, de modo que el tiempo de tránsito baja a unos
\qty{2}{s} --- todavía suficiente para que salga el oxígeno.
\end{solution}

\begin{solution}{exo:b2:blood-circulation:8}
Extremo arterial: $35 - 25 = \qty{+10}{mmHg}$ (filtración); extremo
venoso: $15
- 25 = \qty{-10}{mmHg}$ (reabsorción). Con $P_c = 28$ en el extremo venoso
la presión neta es $+3$: el líquido se filtra a lo largo de todo el
\omterm{def:b2:blood-circulation:vessels}{capilar} y no se reabsorbe nada --- edema, los tobillos hinchados de la
insuficiencia cardíaca.
\end{solution}

\begin{solution}{exo:b2:blood-circulation:9}
$40/66000 = \qty{0.61}{mol/m^3}$; $\pi = 0.61\times 8.314\times 310 =
\qty{1570}{Pa} = \qty{11.8}{mmHg}$; con la mitad de albúmina,
\qty{5.9}{mmHg}. El sodio atraviesa libremente la pared \omterm{def:b2:blood-circulation:vessels}{capilar}, de modo
que su concentración es la misma a ambos lados y no ejerce presión
osmótica a través de ella.
\end{solution}

\begin{solution}{exo:b2:blood-circulation:10}
$RC = \qty{2}{s}$: $120e^{-0.4} = \qty{80}{mmHg}$. $C = 1$: $RC = 1$,
$120e^{-0.8} = \qty{54}{mmHg}$ --- y el mismo volumen sistólico eleva más
la presión sistólica en una aorta rígida. Una \omterm{thm:b2:blood-circulation:windkessel}{presión del pulso} amplia
significa un pico más alto contra el que el corazón debe expulsar la
\omterm{def:b2:blood-circulation:blood}{sangre}, más trabajo y más oxígeno para el mismo gasto.
\end{solution}

\begin{solution}{exo:b2:blood-circulation:11}
Reposo: $90/83 = \qty{1.08}{mmHg.s/mL}$; ejercicio: $105/417 =
\qty{0.25}{mmHg.s/mL}$, una caída de cuatro veces. Como $R \propto r^{-4}$,
$r$ ha aumentado en un factor $4^{1/4} = 1.41$: las \omterm{def:b2:blood-circulation:vessels}{arteriolas} se han
ensanchado en promedio un \qty{41}{\%}.
\end{solution}

\begin{solution}{exo:b2:blood-circulation:12}
La continuidad fija la velocidad de cada nivel mediante la sección total,
y el árbol está construido de modo que la sección es mil veces la de la
aorta en los \omterm{def:b2:blood-circulation:vessels}{capilares} y vuelve a ser pequeña en las \omterm{def:b2:blood-circulation:vessels}{venas}; la \omterm{thm:b2:blood-circulation:poiseuille}{ley de
Poiseuille} sitúa la resistencia en las \omterm{def:b2:blood-circulation:vessels}{arteriolas}, donde el músculo puede
modificarla, y deja casi sin pérdidas los vasos anchos. Una \omterm{def:b2:blood-circulation:circuits}{circulación
abierta} no tiene \omterm{def:b2:blood-circulation:vessels}{capilares} que frenen la \omterm{def:b2:blood-circulation:blood}{sangre} en un lugar definido, ni
vasos en los que fijar una resistencia, ni manera de dirigir el flujo
hacia un órgano --- solo puede remover.
\end{solution}

\begin{solution}{pb:b2:blood-circulation:1}
\textbf{1.} $70\times 72 = \qty{5.0}{L/min}$; \qty{7300}{L} al día.
\textbf{2.} $7300/5 \approx 1450$ veces.
\textbf{3.} Una cuarta parte de \qty{57}{g} a 72 latidos por minuto:
\qty{14}{g}
$\times 4320 = \qty{62}{kg}$ por hora --- aproximadamente el peso de un
hombre.
\textbf{4.} Ningún alimento podría aportar, ni ningún tejido consumir, el
peso de un hombre en \omterm{def:b2:blood-circulation:blood}{sangre} cada hora; la única salida es que la misma
\omterm{def:b2:blood-circulation:blood}{sangre} vuelva al corazón.
\textbf{5.} Un cordón atado alrededor del brazo hincha las \omterm{def:b2:blood-circulation:vessels}{venas} por
debajo, no por encima, de modo que la \omterm{def:b2:blood-circulation:blood}{sangre} de las \omterm{def:b2:blood-circulation:vessels}{venas} se desplaza hacia
el corazón; si se empuja la \omterm{def:b2:blood-circulation:blood}{sangre} de una \omterm{def:b2:blood-circulation:vessels}{vena} hacia la mano, se detiene en
las válvulas y no puede hacerse retroceder --- las válvulas solo permiten
el flujo hacia el corazón.
\textbf{6.} Los \omterm{def:b2:blood-circulation:vessels}{capilares} que unen las \omterm{def:b2:blood-circulation:vessels}{arterias} con las \omterm{def:b2:blood-circulation:vessels}{venas}; Malpighi los
vio en el pulmón de la rana en 1661.
\textbf{7.} Sección $\pi\times 1.25^{2} = \qty{4.9}{cm^2}$; $83/4.9 =
\qty{17}{cm/s}$.
\textbf{8.} $8\times 3\times 10^{-3}\times 0.4\times 8.3\times
10^{-5}/(\pi\times 2.44\times 10^{-8}) = \qty{10}{Pa}$, \qty{0.08}{mmHg}.
\textbf{9.} $R_{1} = 8\times 3\times 10^{-3}\times 10^{-3}/(\pi\times
5.06\times 10^{-20}) = \qty{1.5e14}{Pa.s/m^3}$; en paralelo,
$1.5\times 10^{14}/3\times 10^{6} = \qty{5.0e7}{Pa.s/m^3}$.
\textbf{10.} $8.3\times 10^{-5}\times 5.0\times 10^{7} = \qty{4200}{Pa}
= \qty{31}{mmHg}$.
\textbf{11.} \omterm{def:b2:blood-circulation:vessels}{Capilares} abiertos: $10^{10}$, cada uno de $\pi(3\times 10^{-6})^{2}
= \qty{2.8e-11}{m^2}$: \qty{0.28}{m^2} $= \qty{2800}{cm^2}$; $v =
83/2800 = \qty{0.03}{cm/s}$; tránsito: $0.1/0.03 = \qty{3.3}{s}$.
\textbf{12.} Resistencia $\times(1/0.9)^{4} = 1.52$: \qty{47}{mmHg}; el
corazón debe elevar la presión arterial en \qty{16}{mmHg}, o el gasto cae
un tercio.
\textbf{13.} $+10$, $-10$ y $0$ mmHg.
\textbf{14.} Filtración: $10\times 2 = \qty{20}{L}$ al día; reabsorción:
$8.5\times 2 = \qty{17}{L}$.
\textbf{15.} \qty{3}{L} de linfa al día.
\textbf{16.} $\pi_c = \qty{12.5}{mmHg}$: neta de $+22.5$ en el extremo
arterial y de $+2.5$ en el extremo venoso --- filtración en todas partes,
sin reabsorción: edema.
\textbf{17.} Netas de $+10$ y $+3$, media de $+6.5$: filtración a lo largo
de todo el \omterm{def:b2:blood-circulation:vessels}{capilar}, unos \qty{26}{L} al día; los linfáticos no pueden
transportarlos y el líquido se acumula en los tejidos, empezando por los
pies.
\textbf{18.} $4\times 10^{10}\times 2\pi\times 3\times 10^{-6}\times
10^{-3} = \qty{750}{m^2}$; $t = x^{2}/2D = (2\times 10^{-5})^{2}/
(2\times 10^{-9}) = \qty{0.2}{s}$.
\textbf{19.} $(93 - 3)/83 = \qty{1.08}{mmHg.s/mL}$.
\textbf{20.} $RC = \qty{2.2}{s}$; $120e^{-0.8/2.2} = 120\times 0.69 =
\qty{83}{mmHg}$.
\textbf{21.} $RC = \qty{1.1}{s}$; $120e^{-0.73} = \qty{58}{mmHg}$: la
\omterm{thm:b2:blood-circulation:windkessel}{presión del pulso} se ensancha de 37 a \qty{62}{mmHg}.
\textbf{22.} $120e^{-0.3/2.2} = \qty{105}{mmHg}$: a frecuencia alta la
presión apenas cae entre latidos.
\textbf{23.} $C\Delta P = 2\times 20 = \qty{40}{mL}$ almacenados; los otros
\qty{30}{mL} salen por las \omterm{def:b2:blood-circulation:vessels}{arteriolas} durante la propia sístole.
\textbf{24.} El ventrículo que se contrae aprieta hasta cerrar sus propios
vasos durante la sístole, de modo que el músculo cardíaco se irriga en
diástole, gracias a la presión que mantiene la aorta al retraerse; una
aorta rígida que deja desplomarse la presión diastólica deja al corazón
sin riego entre latidos.
\textbf{25.} Gasto de \qty{5}{L/min}; caída arteriolar de \qty{31}{mmHg};
filtrado neto de \qty{3}{L} al día hacia la linfa; diastólica de
\qty{83}{mmHg} para $C = 2$ y de \qty{58}{mmHg} para $C = 1$.
\end{solution}
